%%% MA_11_Cycles_Crises.tex -- reorganised; frames moved from the old MA_xx files (see comments)

% \section{Business Cycles, Crises, and Stabilisation Policy}

%%% [moved from MA_10_Cycles.tex, frame 0]
\begin{frame}{Business Cycles}
    \mytext{
    \showmytikz[1,2,3]{../shared/tikz/tikz-porkcycle}
    
    Economic growth is not steady, we see phases of expansion and contraction, of boom and bust.

    What exactly determines these cycles? Are they predictable?
    }
    \mycomment{
    In stock broker jargon, markets with rising prices (expanding markets) are nicknamed \enquote{bull markets} -- the bull is a fierce animal \enquote{fighting upwards}. Falling prices (contracting markets) are nicknamed \enquote{bear markets} -- the bear is a fierce animal \enquote{fighting downwards}.

    Longing: you \enquote{long} something (you haven't got it, but you want to have it) -- you are betting on rising prices, you are trying to buy. Long selling: you actually own the asset that you are selling. 

    Shorting: you are trying to sell. Short selling: you are trying to sell without acutally having bought the asset yet -- you should have the option to buy it though (call option). 
    }
\end{frame}

%%% [moved from MA_10_Cycles.tex, frame 1]
\begin{frame}{Anticyclic Policy}
    \mycomment{
    If you know how the cycle runs, you could (in theory) counter-act it, thus evening the amplitude out. If ups and downs are natural and deterministic for lassaiz-faire supply and demand, adding artificial government demand (and supply) could smooth these out. This is the idea of anticyclic fiscal policy.
    }
        \showmytikz[1,2,3]{../shared/tikz/tikz-anticyclic.tex}

\end{frame}

%%% [moved from MA_10_Cycles.tex, frame 2]
\begin{frame}{Kitchin Cycle}
    \mytext{
    \showmytikz[1]{../shared/tikz/tikz-businesscycles}
    
    The Kitchin cycle is an inventory cycle, like the Pork Cycle, with a cycle length of 3-5 years.
    }
\end{frame}

%%% [moved from MA_10_Cycles.tex, frame 3]
\begin{frame}{Juglar Cycle}
    \mytext{
    \showmytikz[2]{../shared/tikz/tikz-businesscycles}
    
    The Juglar cycle is a fixed investment cycle. 

    Cycle length: 7-11 years.
    }

\end{frame}

%%% [moved from MA_10_Cycles.tex, frame 4]
\begin{frame}{Overlapping Cycles}
    \mytext{
    \showmytikz[1,2,3]{../shared/tikz/tikz-businesscycles}
    
    The red curve shows the sum of the cycles already introduced (with a scaled down amplitude).  
    }
\end{frame}

%%% [moved from MA_10_Cycles.tex, frame 5]
\begin{frame}{Kuznets Swing}
    \mytext{
    \showmytikz[4,5]{../shared/tikz/tikz-businesscycles}
    
    Kuznets swings relates to infrastructure investment.

    Cycle length: 15-25 years.
    }
\end{frame}

%%% [moved from MA_10_Cycles.tex, frame 6]
\begin{frame}{Kondratiev Wave}
    \mycomment{
    Kondratiev Waves are larger boom and bust cycles attributed to major technological breakthroughs. The invention of the steam engine, the combustion engine, the computer are much quoted.

    In these discussions, the negative (contractive) phase is often left out. But every new invention also makes existing investment obsolete, causing factories to shut etc. and thus having a negative impact as well.
    }
    \mytext{
    \showmytikz[6,7]{../shared/tikz/tikz-businesscycles}
    
    Kondratiev Waves are (legendary?) long-term ups and downs due to technological change. 

    Cycle length: 45-60 years.
    }
\end{frame}

%%% [moved from MA_10_Cycles.tex, frame 7]
\begin{frame}{Business Cylces: Random Walk}
    \mytext{
    \showmytikz[7,8]{../shared/tikz/tikz-businesscycles}
    In reality, economic activity appears to resemble a random walk (here in green) rather than cycles. If that is so, there is no more argument for a countercyclical policy. 
    }
    \mycomment{
    There is a real difference between \enquote{ordinary} business cycles, and black-swan crisis moments. Random ups and downs might rock the boat. But every now and then, for whatever reasons, waves may sum up creating a perfect storm and capsizing a boat. 

    The Greek word \enquote{crisis}  means the \textit{turning point} of an illness, where the patient either recovers or not.  The illness itself is not a crisis in the strict sense of the word.
    }
\end{frame}

%%% [moved from MA_10_Cycles.tex, frame 8]
\begin{frame}{Business Cycles: Credit Driven}
    \mytext{
    Business cycles could be the result of credit boom and bust, and of the central bank's interest setting behaviour (known as Austrian Business Cycle Theory).

    Then it would be the government \enquote{rocking the boat}, depressing interest rates below their natural level, thus setting wrong price signals (over-optimistic investment, \enquote{malinvestment}). 

    The crisis can then not be fixed by further expansion, but is a necessary correction.
    }
\end{frame}


\begin{frame}{Crisis, what Crisis?}
    \mytext{
    If something cannot go on forever, it will stop. (Herbert \noun{Stein})

    Crises take longer to arrive than you can possibly imagine, but when they do come, they happen faster than you can possibly imagine. (Rüdiger \noun{Dornbusch}; sometimes known as \enquote{Dornbusch's Law}.
    }
\end{frame}

%%% TODO (new content): output gap, recession definition; Minsky / debt deflation; case studies 2008, euro crisis, COVID


\begin{frame}{Crisis: Minsky}
    \mytext{
    Is the economy stable? Is growth stable? Or is the system inherently unstable?

    \begin{itemize}
        \item Minsky: \enquote{stability is destabilising}
        \item seemingly stable growth attracts investment
        \item investment principal and interest due climb exponentially
        \item law of diminishing returns: returns rise first, then fall
        \item every exponential curve overtakes every limited curve at some point
        \item herd behaviour can lead to destabilisation        
    \end{itemize}
    }
\end{frame}


\begin{frame}{Crisis: Minsky}
    \mycomment{
    Consider a growing industry. This will attract finance. As the income increases during the boom, interest rate and debt service can be paid.
    }
    \showmytikz[1]{../shared/tikz/tikz-minsky}
\end{frame}

\begin{frame}{Crisis: Minsky}
    \mycomment{
    At some point, generated income covers interest rate only; to serve debt, new debt must be acquired. 
    }
    \showmytikz[1,2]{../shared/tikz/tikz-minsky}
\end{frame}

\begin{frame}{Crisis: Minsky}
    \mycomment{
    If debt levels, debt service and interest continue to rise exponentially, but income levels fall, debt service and interest can only be financed by ever more debt. This is unsustainable.
    Finally, a crisis occurs. As creditors withdraw their fund, income (that was previously generated as cash flow through more debt) falls (dramatically). 
    
    These phases can be characterised as \enquote{hedge finance}, \enquote{speculative finance}, and \enquote{Ponzi finance}. 
    }
    \showmytikz[1,2,3,4]{../shared/tikz/tikz-minsky}
\end{frame}

%%% TODO (new content): public debt: government budget constraint, debt dynamics (r - g), SGP, Schuldenbremse (2025 reform)

\begin{frame}{AI Bubble}
    \showmytikz{../shared/tikz/tikz-aibubble}
    \mycomment{
    At the time of writing, investment in AI and AI infrastructure is the \enquote{elephant in the room}. What are your thoughts?
    }
\end{frame}